\documentclass[a4paper, serif, 10pt]{moderncv}

%% moderncv
% style and color
\moderncvstyle{banking}
\moderncvcolor{black}

%% page layout/margins
\usepackage[top=2.5cm, bottom=2.5cm, left=1.5cm, right=1.5cm]{geometry}

\nopagenumbers{}

%% Babel config for Indonesian (Bahasa Indonesia) language
% ref:
% - https://latex3.github.io/babel/guides/locale-indonesian.html
\usepackage[indonesian]{babel}

%% fontspec
\usepackage{fontspec}

\IfFontExistsTF{Linux Libertine}{
  \setmainfont[Ligatures={TeX, Common}, Numbers=OldStyle]{Linux Libertine}
}{}
\IfFontExistsTF{Linux Libertine O}{
  \setmainfont[Ligatures={TeX, Common}, Numbers=OldStyle]{Linux Libertine O}
}{}

%% CTeX
% CJK support, used to show CJK characters in the resume
%
% - fontset=none: disable builtin fontset but instead set the CJK font manually
% - heading=false: disable ctex heading
% - punct=kaiming: use kaiming punctuations styles for CJK
% - scheme=plain: use plain scheme, do not override \normalsize font size
% - space=auto: space settings for CJK characters
%
% ref:
% - http://ctan.mirrorcatalogs.com/language/chinese/ctex/ctex.pdf
\usepackage[UTF8, heading=false, punct=kaiming, scheme=plain, space=auto]{ctex}

\IfFontExistsTF{Noto Serif CJK SC}{
  \setCJKmainfont{Noto Serif CJK SC}
}{}
\IfFontExistsTF{Noto Sans CJK SC}{
  \setCJKsansfont{Noto Sans CJK SC}
}{}

%% line spacing
\usepackage{setspace}
\setstretch{1.125}

%% URL styling - use normal text instead of monospace
\urlstyle{same}

% Auto-underline all links
% Defer until \begin{document} so hyperref is already loaded
\AtBeginDocument{%
  \let\oldhref\href
  \renewcommand{\href}[2]{\oldhref{#1}{\underline{#2}}}%
}

%% Patch \httplink and \httpslink to handle full URLs
% The original moderncv macros always prepend http:// or https:// to URLs.
% This causes issues when the URL already has a protocol (e.g., https://example.com)
% resulting in malformed URLs like https://https://example.com.
% This patch checks if the URL already contains "://" and uses it directly if so.
% Note: We use \str_set:Nx (x-expansion) to expand macros like \@homepage first.
\ExplSyntaxOn
\str_new:N \g_yamlresume_url_str

\RenewDocumentCommand{\httpslink}{O{}m}{
  \str_set:Nx \g_yamlresume_url_str {#2}
  \str_if_in:NnTF \g_yamlresume_url_str {://}
    { \url{#2} }
    { \url{https://#2} }
}

\RenewDocumentCommand{\httplink}{O{}m}{
  \str_set:Nx \g_yamlresume_url_str {#2}
  \str_if_in:NnTF \g_yamlresume_url_str {://}
    { \url{#2} }
    { \url{http://#2} }
}
\ExplSyntaxOff

\name{Rizky Aditya Pratama}{}
\title{Software Engineer Berfokus pada Backend, Cloud, dan Arsitektur Skalabel}
\phone[mobile]{+62 812-3456-7890}
\email{rizky.aditya@example.com}
\homepage{https://rizkyaditya.dev}

\address{Jl. Senopati No. 45, Kebayoran Baru, Jakarta Selatan, DKI Jakarta, Indonesia, 12190}{}{}

\extrainfo{{\small \faGithub}\ \href{https://github.com/rizkyadityap}{@rizkyadityap} {} {} {} • {} {} {}
{\small \faLinkedin}\ \href{https://www.linkedin.com/in/rizkyadityapratama}{@rizkyadityapratama}}

\begin{document}

\maketitle

\section{Informasi Dasar}

\cvline{}{\begin{itemize}
\item Software engineer dengan 5+ tahun pengalaman membangun layanan backend yang skalabel dan andal di lingkungan startup serta perusahaan teknologi.
\item Mahir menggunakan Go, Java, dan Node.js, serta berpengalaman menerapkan arsitektur mikroservis, event-driven, dan RESTful API.
\item Berpengalaman mengelola infrastruktur di AWS, Docker, dan Kubernetes, termasuk penerapan CI/CD dan observability.
\item Terbiasa berkolaborasi dalam tim agile, melakukan code review, menulis dokumentasi teknis, dan menjadi mentor bagi engineer junior.
\end{itemize}}
\section{Pendidikan}

\cventry{Agu 2016–Jul 2021}
        {Sarjana, Ilmu Komputer, Nilai: 3.76}
        {Universitas Indonesia}
        {\url{https://www.ui.ac.id/}}
        {}
        {\begin{itemize}
\item Lulus dengan predikat cum laude dan fokus pada mata kuliah rekayasa perangkat lunak, basis data, dan sistem terdistribusi.
\item Menjadi asisten dosen untuk praktikum Pemrograman Dasar dan Struktur Data selama dua semester.
\item Menyelesaikan skripsi berjudul 'Perancangan API Gateway untuk Layanan Mikroservis pada Aplikasi E-Commerce'.
\end{itemize}
\textbf{Mata Kuliah}: Struktur Data dan Algoritma, Basis Data, Sistem Operasi, Jaringan Komputer, Rekayasa Perangkat Lunak, Kecerdasan Buatan, Pemrograman Berorientasi Objek, Matematika Diskret}
\section{Pengalaman Kerja}

\cventry{Jan 2022–Sekarang}
        {Software Engineer}
        {PT Nusantara Digital Solusi}
        {\url{https://www.nusantaradigital.co.id}}
        {}
        {\begin{itemize}
\item Merancang dan mengembangkan RESTful API untuk platform logistik menggunakan Go dan PostgreSQL, melayani 2 juta permintaan per hari.
\item Memigrasikan sistem monolitik menjadi arsitektur mikroservis dengan Kafka dan gRPC, mengurangi waktu deployment hingga 40\%.
\item Menerapkan Circuit Breaker dan Retry Mechanism untuk meningkatkan ketahanan sistem terhadap kegagalan upstream.
\item Berkolaborasi dengan tim DevOps membuat pipeline CI/CD menggunakan GitHub Actions dan Argo CD.
\item Memimpin technical sharing dan menjadi mentor bagi onboarding engineer baru.
\end{itemize}
\textbf{Kata Kunci}: Go, Kubernetes, Kafka, gRPC, PostgreSQL, Microservices, CI/CD, REST API}

\cventry{Agu 2020–Des 2021}
        {Backend Developer}
        {PT Kreatif Media Nusantara}
        {\url{https://kreatifmedia.co.id}}
        {}
        {\begin{itemize}
\item Mengembangkan layanan backend untuk aplikasi media streaming menggunakan Node.js, Express, dan MongoDB.
\item Merancang skema database dan mengoptimalkan query sehingga response time turun dari 800 ms menjadi 200 ms.
\item Mengintegrasikan penyedia pembayaran digital seperti Midtrans dan Xendit untuk fitur langganan premium.
\item Menulis unit test dan integration test menggunakan Jest, meningkatkan code coverage dari 45\% menjadi 85\%.
\item Berpartisipasi aktif dalam sprint planning, daily standup, dan retrospektif bersama tim produk.
\end{itemize}
\textbf{Kata Kunci}: Node.js, Express, MongoDB, Redis, Jest, Payment Gateway, REST API}

\cventry{Jul 2019–Jul 2020}
        {Junior Web Developer}
        {PT Aplikasi Anak Negeri}
        {\url{https://aplikasianaknegeri.co.id}}
        {}
        {\begin{itemize}
\item Membangun fitur frontend dan backend pada sistem informasi desa menggunakan PHP Laravel dan MySQL.
\item Membuat modul manajemen pengguna, inventaris aset, dan pelaporan pendapatan desa.
\item Melakukan debugging serta pemeliharaan aplikasi produksi dan menyusun dokumentasi penggunaan untuk admin desa.
\end{itemize}
\textbf{Kata Kunci}: Laravel, PHP, MySQL, JavaScript, Bootstrap}
\section{Bahasa}

\cvline{Melayu}{Bahasa Ibu / Bilingual \hfill \textbf{Kata Kunci}: Bahasa Indonesia}
\cvline{Inggris}{Tingkat Profesional Penuh \hfill \textbf{Kata Kunci}: TOEFL ITP 600, IELTS 7.0}
\section{Keahlian}

\cvline{Pengembangan Backend}{Ahli \hfill \textbf{Kata Kunci}: Go, Java, Node.js, REST API, gRPC, GraphQL}
\cvline{Arsitektur \& Infrastruktur}{Mahir \hfill \textbf{Kata Kunci}: Docker, Kubernetes, AWS, Terraform, Kafka, Redis}
\cvline{Database}{Mahir \hfill \textbf{Kata Kunci}: PostgreSQL, MySQL, MongoDB, Elasticsearch, Redis}
\cvline{Frontend Web}{Menengah \hfill \textbf{Kata Kunci}: React, TypeScript, Tailwind CSS, Vite}
\section{Penghargaan}

\cventry{Nov 2020}
        {Kementerian Komunikasi dan Informatika RI}
        {Juara 2 Hackathon API Nasional 2020}
        {}
        {}
        {Membangun API inovasi untuk layanan kesehatan berbasis komunitas bersama tim dan menang dari 150+ peserta.}
\section{Sertifikat}

\cventry{Mar 2023}
        {Amazon Web Services}
        {AWS Certified Solutions Architect - Associate}
        {\url{https://aws.amazon.com/certification/}}
        {}
        {}

\cventry{Agu 2022}
        {Google Cloud}
        {Professional Cloud Developer - Google Cloud}
        {\url{https://cloud.google.com/certification}}
        {}
        {}
\section{Publikasi}

\cventry{Jun 2023}
        {Implementasi Service Mesh menggunakan Istio pada Layanan Mikroservis}
        {Jurnal Teknologi Informasi dan Ilmu Komputer}
        {\url{https://jtiik.ub.ac.id/index.php/jtiik}}
        {}
        {Membahas pola penerapan service mesh Istio untuk observability, keamanan, dan traffic management pada layanan mikroservis di lingkungan Kubernetes.}
\section{Referensi}

\cventry{\emaillink[budi.santoso@nusantaradigital.co.id]{budi.santoso@nusantaradigital.co.id}}
        {Engineering Manager di PT Nusantara Digital Solusi}
        {Budi Santoso}
        {+62 811-9876-5432}
        {}
        {Rizky adalah salah satu engineer paling andal di tim. Ia memiliki kemampuan analitis yang kuat dan selalu mengutamakan kualitas kode serta kolaborasi tim.}
\section{Proyek}

\cventry{Mar 2021–Agu 2021}
        {Platform pelaporan keuangan dan aset desa berbasis web}
        {Sirekap Desa}
        {\url{https://github.com/rizkyadityap/sirekap-desa}}
        {}
        {\begin{itemize}
\item Mengembangkan aplikasi berbasis Laravel dan MySQL untuk digitalisasi administrasi desa.
\item Merancang sistem autentikasi multi-level untuk admin desa, camat, dan dinas pemberdayaan masyarakat.
\item Menghasilkan dashboard visualisasi pendapatan dan belanja desa menggunakan Chart.js.
\end{itemize}
\textbf{Kata Kunci}: Laravel, MySQL, Chart.js, Bootstrap}

\cventry{Jan 2022–Jun 2022}
        {Backend microservices untuk aplikasi e-commerce sederhana}
        {GoShop API}
        {\url{https://github.com/rizkyadityap/goshop-api}}
        {}
        {\begin{itemize}
\item Merancang service catalog, cart, order, dan payment menggunakan arsitektur microservices.
\item Menggunakan komunikasi gRPC antar service dan REST API untuk klien.
\item Menerapkan Saga Pattern untuk distributed transaction dan Kafka sebagai event bus.
\end{itemize}
\textbf{Kata Kunci}: Go, gRPC, Kafka, Docker, PostgreSQL}
\section{Minat}

\cvline{Sepak Bola}{Futsal, Premier League}
\cvline{Teknologi Open Source}{Kubernetes, Golang, Linux}
\cvline{Literasi}{Buku Teknologi, Jurnal Ilmiah}
\section{Kegiatan Sukarela}

\cventry{Feb 2022–Des 2023}
        {Relawan Acara}
        {Google Developer Groups Jakarta}
        {\url{https://gdg.community.dev/gdg-jakarta/}}
        {}
        {\begin{itemize}
\item Menjadi relawan untuk Google I/O Extended Jakarta dan Cloud Study Jam.
\item Membantu koordinasi peserta, dokumentasi, serta moderasi sesi diskusi.
\end{itemize}}


\cventry{Jan 2021–Des 2021}
        {Mentor Belajar Backend}
        {Komunitas Programmer Indonesia}
        {\url{https://komunitasprogrammer.id}}
        {}
        {\begin{itemize}
\item Memberikan mentorship gratis bagi anggota komunitas yang ingin mendalami backend development.
\item Membimbing lima mentee dalam menyelesaikan proyek capstone API sederhana.
\end{itemize}}

\end{document}